\documentclass[aspectratio=169]{beamer}

\usepackage[utf8]{inputenc}
\usepackage[T1]{fontenc}
\usepackage{amsmath,amssymb,amsfonts,bm,mathtools}
\usepackage{physics}
\usepackage{braket}
\usepackage{quantikz}
\usepackage{graphicx}
\usepackage{booktabs}
\usepackage{hyperref}

\usetheme{Madrid}

\title{Quantum Simulation and the HHL Algorithm}
\subtitle{with Phase Estimation Circuits and a Worked Example for the 1D Poisson Equation}
\author{}
\date{}

\begin{document}

\frame{\titlepage}

\section{Motivation}

\begin{frame}{Feynman's Vision}
For an \(n\)-qubit system,
\[
\ket{\psi}=\sum_{z\in\{0,1\}^n} \alpha_z \ket{z},
\qquad
\sum_z |\alpha_z|^2=1,
\]
so one must track \(2^n\) complex amplitudes.

\medskip

This motivates \emph{quantum simulation}: use controllable quantum systems to simulate other quantum systems.
\end{frame}

\begin{frame}{Schrödinger Dynamics and Hamiltonian Simulation}
The dynamics are governed by
\[
i\frac{d}{dt}\ket{\psi(t)} = H\ket{\psi(t)},
\]
with formal solution
\[
\ket{\psi(t)} = e^{-iHt}\ket{\psi(0)}.
\]

Hamiltonian simulation asks for a quantum circuit implementing
\[
U(t)\approx e^{-iHt}
\]
efficiently in the number of qubits, simulation time, and target precision.
\end{frame}

\begin{frame}{From Simulation to Linear Systems}
Many physics and numerical analysis problems reduce to
\[
A\bm{x}=\bm{b}.
\]

Examples include:
\begin{itemize}
    \item finite-difference discretizations of Poisson equations,
    \item implicit time stepping for diffusion equations,
    \item linear response theory.
\end{itemize}

The Harrow--Hassidim--Lloyd (HHL) algorithm prepares
\[
\ket{x}\propto A^{-1}\ket{b}.
\]
\end{frame}

\section{Quantum Phase Estimation}

\begin{frame}{Phase Estimation Setup}
Suppose \(U\) is unitary and
\[
U\ket{u}=e^{2\pi i \phi}\ket{u},
\qquad
\phi\in[0,1).
\]

Quantum phase estimation estimates \(\phi\) using:
\begin{itemize}
    \item a control register of \(m\) qubits,
    \item controlled powers \(U^{2^k}\),
    \item the inverse quantum Fourier transform.
\end{itemize}
\end{frame}

\begin{frame}{Phase Estimation Circuit}
\[
\begin{quantikz}[row sep=0.35cm, column sep=0.45cm]
\lstick{\ket{0}} & \gate{H} & \ctrl{3} & \qw      & \qw      & \gate[wires=3]{QFT^\dagger} & \meter{} \\
\lstick{\ket{0}} & \gate{H} & \qw      & \ctrl{2} & \qw      &                              & \meter{} \\
\lstick{\ket{0}} & \gate{H} & \qw      & \qw      & \ctrl{1} &                              & \meter{} \\
\lstick{\ket{u}} & \qw      & \gate{U} & \gate{U^2} & \gate{U^4} & \qw                        & \qw
\end{quantikz}
\]
\end{frame}

\begin{frame}{How the Phases Accumulate}
After the Hadamard gates, the control register is
\[
\frac{1}{2^{m/2}} \sum_{k=0}^{2^m-1} \ket{k}.
\]

Applying the controlled powers gives
\[
\frac{1}{2^{m/2}}\sum_{k=0}^{2^m-1} e^{2\pi i k\phi}\ket{k}\ket{u}.
\]

The inverse QFT maps this state approximately to
\[
\ket{\widetilde{\phi}}\ket{u}.
\]
\end{frame}

\section{The HHL Algorithm}

\begin{frame}{The Quantum Linear Systems Problem}
Let
\[
A\ket{x}=\ket{b},
\]
where \(A\) is Hermitian and invertible.

Write
\[
A=\sum_j \lambda_j \ket{u_j}\bra{u_j},
\qquad
\ket{b}=\sum_j \beta_j \ket{u_j}.
\]

Then
\[
\ket{x}\propto A^{-1}\ket{b}
=
\sum_j \frac{\beta_j}{\lambda_j}\ket{u_j}.
\]
\end{frame}

\begin{frame}{Main Idea Behind HHL}
The essential map is
\[
\lambda_j \mapsto \frac{1}{\lambda_j}.
\]

Since this is not unitary, HHL proceeds indirectly:
\begin{enumerate}
    \item perform phase estimation on \(e^{iAt}\),
    \item store an approximation of \(\lambda_j\),
    \item perform a controlled rotation whose amplitude depends on \(1/\lambda_j\),
    \item uncompute the phase register,
    \item postselect on a successful ancilla measurement.
\end{enumerate}
\end{frame}

\begin{frame}{Full HHL Circuit Schematic}
\[
\begin{quantikz}[row sep=0.35cm, column sep=0.38cm]
\lstick{\ket{0}} & \qw & \gate[style={fill=blue!10}]{R(\lambda^{-1})} & \qw & \meter{} \\
\lstick{\ket{0}^{\otimes m}} & \gate[wires=2, style={fill=green!10}]{\mathrm{QPE}(e^{iAt})} & \qw & \gate[wires=2, style={fill=green!10}]{\mathrm{QPE}^\dagger} & \qw \\
\lstick{\ket{b}} & \qw & \qw & \qw & \rstick{\ket{x}\propto A^{-1}\ket{b}}
\end{quantikz}
\]
\end{frame}

\begin{frame}{HHL State Evolution}
Start with
\[
\ket{0}\ket{0}^{\otimes m}\ket{b}
=
\ket{0}\ket{0}^{\otimes m}\sum_j \beta_j \ket{u_j}.
\]

After phase estimation:
\[
\ket{0}\sum_j \beta_j \ket{\widetilde{\lambda_j}}\ket{u_j}.
\]

After controlled rotation:
\[
\sum_j \beta_j
\left(
\sqrt{1-\frac{C^2}{\lambda_j^2}}\ket{0}
+
\frac{C}{\lambda_j}\ket{1}
\right)
\ket{\widetilde{\lambda_j}}\ket{u_j}.
\]

After uncomputation and postselection on ancilla \(\ket{1}\):
\[
\ket{x}\propto \sum_j \frac{\beta_j}{\lambda_j}\ket{u_j}.
\]
\end{frame}

\section{Worked Example: 1D Poisson Equation}

\begin{frame}{Poisson Equation in One Dimension}
Consider
\[
- u''(x)=f(x),
\qquad
x\in(0,1),
\]
with Dirichlet boundary conditions
\[
u(0)=u(1)=0.
\]

A simple benchmark choice is
\[
f(x)=\pi^2 \sin(\pi x),
\]
whose exact solution is
\[
u(x)=\sin(\pi x).
\]
\end{frame}

\begin{frame}{Finite-Difference Discretization}
Choose interior grid points
\[
x_i = ih,
\qquad
h=\frac{1}{N+1},
\qquad
i=1,\dots,N.
\]

Use the centered second derivative:
\[
u''(x_i)\approx \frac{u_{i+1}-2u_i+u_{i-1}}{h^2}.
\]

Then
\[
-\frac{u_{i+1}-2u_i+u_{i-1}}{h^2}=f_i,
\]
which yields
\[
A\bm{u}=\bm{b},
\]
with
\[
A=\frac{1}{h^2}
\begin{pmatrix}
2 & -1 & 0 & \cdots & 0\\
-1 & 2 & -1 & \ddots & \vdots\\
0 & -1 & 2 & \ddots & 0\\
\vdots & \ddots & \ddots & \ddots & -1\\
0 & \cdots & 0 & -1 & 2
\end{pmatrix}.
\]
\end{frame}

\begin{frame}{A Small HHL-Compatible Example: \(N=2\)}
Take \(N=2\). Then
\[
h=\frac{1}{3},
\qquad
A = 9
\begin{pmatrix}
2 & -1\\
-1 & 2
\end{pmatrix}
=
\begin{pmatrix}
18 & -9\\
-9 & 18
\end{pmatrix}.
\]

For \(f(x)=\pi^2\sin(\pi x)\),
\[
\bm{b}
=
\pi^2 \frac{\sqrt{3}}{2}
\begin{pmatrix}
1\\
1
\end{pmatrix}.
\]
\end{frame}

\begin{frame}{Exact Classical Solution for \(N=2\)}
The matrix
\[
A=
\begin{pmatrix}
18 & -9\\
-9 & 18
\end{pmatrix}
\]
has eigenpairs
\[
\lambda_1=9,
\qquad
\ket{u_1}=\frac{1}{\sqrt{2}}
\begin{pmatrix}
1\\
1
\end{pmatrix},
\]
\[
\lambda_2=27,
\qquad
\ket{u_2}=\frac{1}{\sqrt{2}}
\begin{pmatrix}
1\\
-1
\end{pmatrix}.
\]

Since
\[
\bm{b}\propto
\begin{pmatrix}
1\\1
\end{pmatrix},
\]
we have
\[
A^{-1}\bm{b}\propto \frac{1}{9}\ket{u_1}.
\]

Hence the normalized solution state is
\[
\ket{x}=\ket{u_1}
=
\frac{1}{\sqrt{2}}
\begin{pmatrix}
1\\1
\end{pmatrix}.
\]
\end{frame}

\begin{frame}{Encoding the Poisson Problem on One Qubit}
For \(N=2\), the solution vector has two components, so it fits in one system qubit:
\[
\ket{0}\leftrightarrow
\begin{pmatrix}
1\\0
\end{pmatrix},
\qquad
\ket{1}\leftrightarrow
\begin{pmatrix}
0\\1
\end{pmatrix}.
\]

The normalized right-hand side is
\[
\ket{b} = \frac{1}{\sqrt{2}}(\ket{0}+\ket{1}) = \ket{+},
\]
and the normalized exact solution state is also
\[
\ket{x}=\ket{+}.
\]
\end{frame}

\begin{frame}{Pauli Representation of the Matrix}
The matrix
\[
A=
\begin{pmatrix}
18 & -9\\
-9 & 18
\end{pmatrix}
\]
can be written as
\[
A = 18 I - 9 X.
\]

Hence
\[
e^{iAt}=e^{i(18I-9X)t}=e^{i18t}e^{-i9Xt}.
\]

The global phase \(e^{i18t}\) is irrelevant; the nontrivial evolution is generated by \(X\).
\end{frame}

\begin{frame}{Choosing the Evolution Time}
Choose
\[
t=\frac{\pi}{18}.
\]

Then
\[
\lambda_1 t = \frac{\pi}{2},
\qquad
\lambda_2 t = \frac{3\pi}{2},
\]
so the corresponding eigenphases are well separated modulo \(2\pi\), making phase estimation straightforward in this toy model.
\end{frame}

\begin{frame}{Worked HHL Logic for the Poisson Example}
Because \(\ket{b}=\ket{u_1}\) already lies in a single eigenstate of \(A\), the HHL workflow simplifies:

\begin{enumerate}
    \item Prepare \(\ket{b}=\ket{+}\).
    \item Phase estimation on \(e^{iAt}\) returns \(\lambda_1=9\).
    \item Controlled rotation multiplies the successful ancilla amplitude by \(1/\lambda_1=1/9\).
    \item Uncomputation removes the phase register.
    \item Postselection returns the state \(\ket{u_1}=\ket{+}\).
\end{enumerate}

The overall factor \(1/9\) is absorbed into normalization and the success probability.
\end{frame}

\begin{frame}{What Observable Can We Extract?}
Since HHL returns a state rather than all entries of \(\bm{u}\), one often computes observables.

For
\[
\ket{x}=\ket{+},
\]
we have
\[
\expval{X}{x}=1,
\qquad
\expval{Z}{x}=0.
\]

These match the exact normalized classical solution vector
\[
\frac{1}{\sqrt{2}}
\begin{pmatrix}
1\\
1
\end{pmatrix}.
\]
\end{frame}

\begin{frame}{Summary}
\begin{itemize}
    \item Hamiltonian simulation implements \(e^{-iHt}\).
    \item Quantum phase estimation extracts eigenphases of a unitary.
    \item HHL combines Hamiltonian simulation, phase estimation, controlled rotation, uncomputation, and postselection.
    \item A discretized 1D Poisson equation gives a concrete and instructive worked example.
\end{itemize}
\end{frame}

\begin{frame}{Take-Home Message}
\[
\boxed{
\text{HHL is best understood as a spectral-filter algorithm for } A^{-1}.
}
\]

For differential equations, discretization turns the PDE into a sparse linear system, and HHL then produces a quantum state proportional to the numerical solution.
\end{frame}

\end{document}
